\section*{Chapitre \ref{ch:m3:metals} --- Métaux}

\begin{solution}{exo:m3:metals:1}
Le tom tombe le vendredi (le jour ouvré suivant) et le cash le lundi suivant (deux jours ouvrés).
\end{solution}

\begin{solution}{exo:m3:metals:2}
Il vend 500 tonnes à tom (livraison des warrants le lendemain contre cash) et achète 500 tonnes pour le
prompt suivant, le jour d'après (réception en retour)~: il prête du métal pendant un jour et touche la
backwardation, ou paie le contango, entre les deux dates.
\end{solution}

\begin{solution}{exo:m3:metals:3}
Le métal annulé sort du stock livrable~: il reste moins de warrants pour que les vendeurs à découvert livrent aux
échéances proches, ils doivent racheter ou emprunter, ce qui pousse les prix proches au-dessus des prix plus lointains (backwardation).
\end{solution}

\begin{solution}{exo:m3:metals:4}
Lundi~: $2\,000 \times (48\,078 - 28\,919) = $ 38.32~millions d'USD~; au sommet de mardi, un supplément de $2\,000
\times 53\,287 = $ 106.57~millions d'USD.
\end{solution}

\begin{solution}{exo:m3:metals:5}
Vendredi +6.8~\%, lundi +66.3~\% (clôture à clôture), mardi jusqu'au sommet +110.8~\%~: le sommet valait 3.74 fois
le cours d'ouverture de vendredi, soit une hausse de 274~\%.
\end{solution}

\begin{solution}{exo:m3:metals:6}
L'or de Londres était en lingots de 400 onces que le COMEX n'accepte pas~; les refondre exigeait des
raffineries suisses qui avaient fermé, et les transporter demandait des vols qui étaient cloués au sol. Sans
moyen de livrer, l'écart n'était pas un arbitrage mais un pari sur la durée de la perturbation, avec des
appels de marge sur les contrats vendus à découvert dans l'intervalle.
\end{solution}

\begin{solution}{exo:m3:metals:7}
65 dates prompt quotidiennes, 14 hebdomadaires et 21 mensuelles (week-ends seuls, jours fériés non pris en compte).
\end{solution}

\begin{solution}{exo:m3:metals:8}
Couvert en valeur, pas en cash. Ses contrats à terme vendus appelaient une marge de variation immédiate, alors que la hausse de la
valeur de sa production future ne serait encaissée qu'à la vente du métal. Sans
lignes de crédit confirmées pour financer la marge, un producteur couvert peut être contraint de clôturer sa couverture, ou de faire défaut, au
pire cours.
\end{solution}

\begin{solution}{pb:m3:metals:1}
\textbf{1.} 289.19~millions d'USD. \textbf{2.} 18.39~millions d'USD, de l'ouverture à la clôture du vendredi.
\textbf{3.} 191.59~millions d'USD. \textbf{4.} 532.87~millions d'USD. \textbf{5.} 742.85~millions d'USD.
\textbf{6.} Environ 2.6~milliards d'USD le 4 mars et 7.05~milliards d'USD le 7 mars. \textbf{7.} Environ
19.75~milliards d'USD. \textbf{8.} Des transactions étaient conclues au-delà des limites de variation de cours, et le personnel les jugeait
authentiques, et non des erreurs. \textbf{9.} La suspension arrêtait les transactions suivantes, mais la marge aux cours de la nuit
serait restée exigible et aurait pu provoquer plusieurs défauts de membres~; l'annulation permettait de fixer
la marge sur la clôture du 7 mars. \textbf{10.} Acheteurs ayant échangé pendant la nuit, parmi lesquels des sociétés qui
ont engagé des poursuites~; la High Court (2023) et la Court of Appeal (2024) ont jugé les décisions légales.
\textbf{11.} Un vendeur à découvert contraint d'acheter pour couvrir une marge ou une livraison est un acheteur forcé~: d'autres achètent avant lui
et lui vendent plus cher. \textbf{12.} Elle oblige les positions longues dominantes à prêter à une
backwardation plafonnée~; le vendeur était ici en grande partie de gré à gré et le short squeeze portait sur le
\omterm{def:m3:metals:prompt}{prix à trois mois}, et non seulement sur le spread proche. \textbf{13.} Ces positions ne sont pas visibles de la bourse, et
les données de positions de la bourse ne pouvaient pas montrer la totalité de la position vendeuse. \textbf{14.} Une liquidité faible et un personnel junior
laissaient le cours s'emballer et les limites être désactivées. \textbf{15.} Un plan de soutien des
banques du producteur. \textbf{16.} Lignes de crédit confirmées ou cash pour la marge, ou options plutôt que contrats à terme pour une partie
de la couverture. \textbf{17.} Personnel senior formé, présent à toute heure, limites automatiques de variation de cours qui ne puissent pas être
relevées pour accompagner une envolée des cours, et une surveillance renforcée des expositions de gré à gré.
\textbf{18.} Seulement pour empêcher un marché désordonné de provoquer des défauts, et selon des règles connues à
l'avance~: l'annulation est elle-même un risque que chaque trader doit intégrer dans ses prix. \textbf{19.} \emph{Résultat à nommer~:}
532.87~millions d'USD de la clôture du lundi au sommet, et 724.46~millions d'USD de la clôture du vendredi.
\textbf{20.} Marge en cash sur une couverture saine en valeur.
\end{solution}

\begin{solution}{iq:m3:metals:1}
Il s'agit d'un forward sur une date prompt donnée (quotidienne jusqu'à trois mois, puis hebdomadaire, puis mensuelle), réglé
par la livraison de warrants à la date prompt, les profits et pertes étant réalisés à la date prompt plutôt que
quotidiennement~; un contrat à terme CME est un contrat mensuel standard réévalué au marché chaque jour.
\iqlookfor{la structure des dates et le règlement-livraison.}
\end{solution}

\begin{solution}{iq:m3:metals:2}
Les vendeurs à découvert contraints d'acheter (marge, livraison) font face à d'autres acheteurs qui achètent avant eux. Nickel a montré
que les acheteurs forcés peuvent être des hedgers, que la position vendeuse peut rester hors bourse sans que personne la voie, et
qu'une bourse peut y mettre fin en annulant des transactions.
\iqlookfor{achats forcés, positions cachées, intervention de la bourse.}
\end{solution}

\begin{solution}{iq:m3:metals:3}
Les spreads proches glissent vers la backwardation à mesure que le métal livrable se réduit. Transaction~: acheter la date proche
contre une date plus lointaine (position d'emprunt), dimensionnée selon les règles de prêt et le risque que le métal
retourne sous forme de \omterm{def:m3:metals:warrant}{warrant d'entrepôt}.
\iqlookfor{stocks et spreads~; le plafond de la règle de prêt.}
\end{solution}

\begin{solution}{iq:m3:metals:4}
Scénarios de fortes hausses de cours sur plusieurs jours (plusieurs fois le pire historique), marge quotidienne avec
appels intrajournaliers, cash disponible et lignes confirmées, le délai pour lever davantage de fonds, et le seuil auquel
la couverture doit être réduite~; rendre compte de l'horizon de survie.
\iqlookfor{cash, non la valeur~; trajectoires sur plusieurs jours~; sources de financement.}
\end{solution}

\begin{solution}{iq:m3:metals:5}
La livraison au COMEX exigeait des lingots de 100 onces, Londres en détient de 400 onces, et le raffinage et les vols étaient
perturbés, si bien que l'arbitrage était impossible. L'écart s'est résorbé quand la logistique a repris et qu'un nouveau contrat
a accepté les lingots de 400 onces et les lingots d'un kilogramme.
\iqlookfor{le lien physique derrière un spread financier.}
\end{solution}

\begin{solution}{iq:m3:metals:6}
Mouvements de cours très en dehors des distributions historiques, dépassements de limites, flux d'ordres à sens unique, insuffisances de marge des
membres de la veille, positions vendeuses concentrées déclarées par les membres, et
transactions aux heures creuses. Automatiser les alertes sur ces signaux et leur remontée au personnel senior, et conserver des limites strictes
que le personnel de nuit ne puisse pas désactiver.
\iqlookfor{signaux concrets, escalade et limites strictes.}
\end{solution}
